\section{Affine decoders in a pretrained model track answer sharing, not exposure}
\label{sec:lment-frequency}

\citet{merullo2025frequency} report that relations with more frequent
subject--object co-occurrences have better linear relational embeddings (LREs).
Writing their OLMo-7B relation-level comparison as
\begin{equation}
  \begin{gathered}
    \bar f_r=\frac{1}{|\mathcal D_r|}\sum_{(s,o)\in\mathcal D_r}n(s,o), \\
    \operatorname{corr}_r(\log\bar f_r,C_r)\approx0.82,
  \end{gathered}
  \label{eq:merullo-frequency}
\end{equation}
$n(s,o)$ counts co-occurrences and $C_r$ measures how successfully LRE-based
interventions change the predicted object. These observations suggest exposure
frequency as an explanation for relational linearity.

We examine this in LMEnt-1B, a one-billion-parameter model pretrained on English
Wikipedia \citep{gottesman2025lment}. We evaluate 5,388 facts across 58 relations
using mean-Jacobian affine decoders $\hat o_r(s)=D(\beta W_r h_s+b_r)$
\citep{hernandez2024linearity}, where $D$ decodes the top-1 token. Our score is
first-token faithfulness, $\phi_r=\Pr[\hat o_r(s)=o_{\mathrm{LM}}(s,r)]$;
frequency counts subject-linked corpus documents containing the gold object label.

\paragraph{The score needs a baseline.}
$\phi_r$ is uncorrected for a decoder that always emits a relation's most common
answer, and only 14 of 30 relations with at least 30 facts beat that constant
guess. Person employer and astronomical parent body reach $\phi_r=1.000$ because
the model emits one answer token for every fact in them. We therefore screen
relations on answer and decoder diversity alone, never on corpus counts: we keep a
relation with at least 30 facts, five distinct model answers, ten distinct decoder
outputs, and faithfulness at least five points above its modal baseline. Eleven of
58 relations survive, and ten support the per-fact analyses below, covering
1,826 facts.

\paragraph{Frequency is a weak relation-level predictor.}
Causal success correlates with $\log(1+\bar f_r)$ at $0.393$ over 56 relations and
15{,}874 held-out trials (Figure~\ref{fig:lment-frequency}A), and faithfulness at
$0.095$ over 48 relations with at least two distinct predicted answers.
Within-relation frequency quartiles move in both directions
(Figure~\ref{fig:lment-frequency}B); those comparisons are not corrected for the
modal baseline and we do not rely on them.

\paragraph{What predicts decoder success is how many entities share an answer.}
For each fact we count the other entities in its relation that receive the same
model answer. Faithfulness rises monotonically with that count
(Table~\ref{tab:lment-sharing}). In a logistic regression over the 1{,}219
non-modal facts, with relation fixed effects and standardized predictors, answer
sharing gives $+0.598$ (95\% CI $[+0.398,+0.796]$) while subject--object
co-occurrence gives $+0.026$ (95\% CI $[-0.159,+0.226]$). Holding answer sharing
fixed, exposure is indistinguishable from zero. A decoder succeeds when an answer
is a category spanning many entities, not when a pair was seen often.

\begin{table}[t]
\centering
\small
\begin{tabular}{lrr}
\toprule
Other entities sharing the answer & Facts & $\phi$ \\
\midrule
0 (answer unique to this entity) & 176 & 0.409 \\
1--9   & 422 & 0.524 \\
10--29 & 357 & 0.692 \\
30+    & 871 & 0.762 \\
\bottomrule
\end{tabular}
\caption{Decoder faithfulness by how many other entities in the same relation
receive the same model answer, over 1,826 facts in ten screened relations.}
\label{tab:lment-sharing}
\end{table}

\paragraph{Both synthetic conditions occur inside one relation.}
Lavertezzo is a municipality in the Italian-speaking canton of Ticino. It has 13
linked corpus documents and four containing the object label; an LLM reading all 13
finds none that states its official language and nine that place it in Ticino. Both
the model and the affine decoder answer \texttt{Italian}. Switzerland, in the same
relation and under the same operator, has 24{,}928 subject--object co-occurrences
across 106{,}201 documents. The model answers \texttt{French} correctly, the
decoder answers \texttt{Swiss}, and a linear probe on the same subject state also
fails. Low exposure with shared structure is decodable, and heavy exposure without
it is not. The Ticino result is not in-context interpolation: rebuilt from
demonstrations containing no Swiss entity and no Italian example, the decoder still
returns \texttt{Italian} for all six held-out Ticino municipalities.

\paragraph{The decoder measures separability, and measures it off-axis.}
A leave-one-out nearest-centroid classifier on the same subject states scores
$0.581$ across the ten screened relations, against $0.575$ for the affine decoder:
the relation-specific operator adds little over how well subject representations
already cluster by answer. It is also misaligned. Pulling each answer contrast back
through the decoder's own mean Jacobian and comparing it with the corresponding
centroid difference gives a mean principal-angle cosine of $0.341$ and a subspace
overlap of $0.136$, against $0.022$ for random directions in 2048 dimensions.
Decoder failure is therefore weak evidence about representation: of 622 facts the
decoder misses, a probe recovers 282 ($45\%$). Claims that a relation is not
affinely represented should be made against a probe rather than against the decoder
alone.

These measurements identify which property tracks affine decodability in a
pretrained model, but cannot manipulate it, since exposure, answer sharing and
entity prominence are confounded in natural text. The controlled corpora that
follow vary structure and exposure independently.
Appendix~\ref{app:lment-frequency} gives the evaluation and aggregation details.

\begin{figure*}[t]
  \centering
  \includegraphics[width=\textwidth]{lment_frequency_aggregate.pdf}
  \caption{Frequency, LRE causality, and faithfulness in LMEnt-1B.
  \textbf{A:} Mean subject--object document count versus pairwise intervention
  success for 56 relations in the held-out causal evaluation; each relation has
  equal weight in the correlation. The score is on $[0,1]$ and does not measure
  exact full-object generation.
  \textbf{B:} Upper-minus-lower-quartile faithfulness for every eligible relation
  in the historical evaluation on three frequency axes. Each tail has at least ten facts. Diamonds show means;
  counts above columns give decreases, ties, and increases. These quartile
  comparisons are not corrected for the modal-answer baseline.}
  \label{fig:lment-frequency}
\end{figure*}
